% Revised 2026-11-11: integrated Results, prose table introductions, and adjacent Tables 4--5.
\subsection{\rev{Automatic calibration and coefficient transfer}}
\label{sec:calibration_results}
\begin{revisionblock}
To assess how soft support depends on its coefficients, we compare equal, learned, and leave-one-dataset-out (LODO) weights for a common five-feature scorer. Learned weights are fitted on validation data with query-group cross-validation; LODO uses the other four datasets without target calibration labels. Candidate membership remains fixed. \tabref{tab:calibration_main} summarizes candidate MRR on the complete test splits alongside the original support formula.

% Keep the two summary tables together, without intervening prose.
\end{revisionblock}
\begin{tableorg}[!t]
\begin{revisionblock}
\par\noindent\begin{minipage}{\linewidth}

\centering\small
\captionof{table}{\textbf{Candidate calibration on the complete test splits.} MRR@20 is computed over all rows, with zero reciprocal rank for absent gold entities. Shared equal, learned, and LODO use the same features; comparison with the original soft-support formula also changes features. LODO is coefficient transfer without target calibration labels.}
\label{tab:calibration_main}
\setlength{\tabcolsep}{6pt}
\begin{tabular}{lrrrrr}
\toprule
Dataset & Gold@20 & Original soft & Shared equal & Learned & LODO \\
\midrule
PrimeKG & 0.240 & 0.12533 & 0.12519 & 0.12533 & \textbf{0.12610} \\
PharmKG & 0.092 & 0.02816 & \textbf{0.03188} & \textbf{0.03188} & 0.03186 \\
Hetionet & 0.440 & 0.09544 & 0.10127 & \textbf{0.11528} & 0.11358 \\
DRKG & 0.200 & 0.04656 & \textbf{0.06143} & 0.06050 & 0.05861 \\
repoDB & 0.214 & 0.04811 & 0.05737 & \textbf{0.06490} & \textbf{0.06490} \\
\bottomrule
\end{tabular}


\end{minipage}\par\medskip
\par\noindent\begin{minipage}{\linewidth}

\centering\small
\captionof{table}{\textbf{Three-hop reachability at matched unique-edge budgets.} Reachability is the per-row fraction of source-reachable candidates still reachable within three edges. $N$ counts eligible rows; absent source connectivity is excluded. Random is mean $\pm$ standard deviation over five seeds. Fuzzy candidate-incidence retention relative to the source is 1.0000 on every dataset.}
\label{tab:retention_main}
\setlength{\tabcolsep}{5pt}
\begin{tabular}{lrrrrr}
\toprule
Dataset & $N$ & Fuzzy & Random & Shortest paths & Shortest coverage \\
\midrule
PrimeKG & 488 & 0.7737 & $0.7760\pm0.0038$ & \textbf{0.9985} & 0.9993 \\
PharmKG & 500 & 0.9010 & $0.4423\pm0.0064$ & \textbf{1.0000} & 1.0000 \\
Hetionet & 90 & 0.4835 & $0.4848\pm0.0428$ & \textbf{1.0000} & 0.8514 \\
DRKG & 386 & 0.8921 & $0.5043\pm0.0170$ & \textbf{1.0000} & 0.8199 \\
repoDB & 297 & 0.9546 & $0.5465\pm0.0241$ & \textbf{1.0000} & 0.8971 \\
\bottomrule
\end{tabular}


\end{minipage}\par\medskip
\end{revisionblock}
\end{tableorg}
\begin{revisionblock}

Shared learned calibration improves candidate MRR over the original formula on four datasets and matches it on PrimeKG; LODO yields higher point estimates on all five. Equal weights already improve four datasets, so these comparisons reflect both the shared features and coefficient learning. Relative to equal weights, the learned gain on repoDB is 0.00754, with an exploratory query-cluster 95\% interval of [0.00040, 0.01518]; the other four intervals include zero.

We also evaluate validation-based selection among 12 edge-coefficient vectors at fixed candidate order and graph budget. Supplementary Section S12 and Tables S16--S24 detail the fitting procedures, coefficient transfer, sensitivity, and downstream comparisons. Together, these analyses provide alternatives to manual weight assignment within the prepared benchmark pipelines.
\end{revisionblock}

\subsection{\rev{Information retention and alternative graph selection}}
\label{sec:retention_results}
\begin{revisionblock}
To examine what fuzzy retrieval retains as the graph becomes smaller, we compare alternative selectors at matched per-query unique-edge budgets with candidate identities and order fixed. The alternatives include random and coverage-constrained sampling, native-score and degree selection, breadth-first search, personalized PageRank, shortest paths, and relation-stratified selection. \tabref{tab:retention_main} reports three-hop query--candidate reachability and highlights the trade-off with candidate-incidence coverage.

Fuzzy retrieval retains all source-covered candidate incidence and more three-hop reachability than uniform random selection on PharmKG, DRKG, and repoDB; PrimeKG and Hetionet are similar to random. Shortest-path selection preserves nearly all source connectivity but reduces candidate-incidence coverage on Hetionet, DRKG, and repoDB. Thus, retaining evidence for every candidate and preserving connected paths to the query favor different selection policies.

Supplementary Section S13, Tables S25--S29, and Figs.~S4--S6 extend this analysis to labeled paths, encoder representations, downstream behavior, and graph budgets. They also distinguish duplicate removal from unique-triple removal. These comparisons characterize fuzzy retrieval as coverage-oriented evidence compression whose structural retention depends on graph topology.
\end{revisionblock}
